\section{The Topology-Aware QuaRK Family}
\label{sec:quark-family}

A QuaRK run separates representational choices from learning budgets.  The
\emph{architecture}
\[
    \arch:=(n,R,\lambda_+),
\]
collects the qubit width, number of parallel reservoir branches, and upper
memory endpoint.  The measurement repetitions $M$, frozen-pool size $S$, and
Ivanov radius $\Lambda$ remain external budgets: they control estimation,
random search, and classical readout capacity rather than the frozen quantum
representation itself.  This distinction will keep the resource dependence of
the later guarantees explicit without repeatedly expanding its components.

\subsection{Gaussian preprocessing and interaction graph}

To compare widths without changing the marginal model at any fixed $n$, draw
an i.i.d. sequence $\vg_1,\vg_2,\ldots\sim\gN(0,I_d)$, independently of the
temporal process, and use the canonical Gaussian prefixes
\begin{equation}
    \projection_n
    :=\frac1{\sqrt n}
      \begin{bmatrix}
        \vg_1^\top\\[-1mm]
        \vdots\\[-1mm]
        \vg_n^\top
      \end{bmatrix}
      \in\sR^{n\times d},
    \qquad n\geq3.
    \label{eq:gaussian-projection}
\end{equation}
For every fixed width, the entries are i.i.d. $\gN(0,1/n)$, exactly as in a
direct Gaussian draw.  The coupling only gives a coherent meaning to
increasing width: the narrower representation is recoverable from the wider
one through
\begin{equation}
    \projection_n x
      =\sqrt{\frac{n+1}{n}}\,
        [\projection_{n+1}x]_{1:n}.
    \label{eq:projection-prefix-recovery}
\end{equation}
A run freezes the prefix required by its chosen width.  With the injective
bounded angle map $\vartheta(u):=\pi\tanh u$ and
$\mR_y(\theta):=\exp(-\ii\theta\pauliY/2)$, encode $x\in\gX$ by
\begin{equation}
    \inputUnitary_n(x)
    :=\bigotimes_{j=1}^n
      \mR_y\!\left(\vartheta((\projection_nx)_j)\right).
    \label{eq:input-encoding}
\end{equation}

The interaction graph is the cycle $\gG_n=(\gV_n,\gE_n)=C_n$, with
$\gV_n=\{1,\ldots,n\}$.  A minimum edge colouring writes
\[
    \gE_n=\bigsqcup_{j=1}^{J_n}\gM_{n,j},
    \qquad
    J_n=\chi'(C_n)=
    \begin{cases}
        2,&n\text{ even},\\
        3,&n\text{ odd},
    \end{cases}
\]
where each $\gM_{n,j}$ is a matching.  The same cycle determines both the
entangling layers and the two-qubit observable bank.

\subsection{One balanced random mixing block}
\label{sec:single-balanced-circuit}

Let $\pauliVec=(\pauliX,\pauliY,\pauliZ)$ and
$\sS^2:=\{\vu\in\sR^3:\norm{\vu}_2=1\}$.  For a vertex $v$ and edge
$e=\{v,v'\}$, use the local generators
$\mG_v(\va):=\va\cdot\pauliVec_v$ and
$\mG_e(\va,\va'):=(\va\cdot\pauliVec_v)\otimes
(\va'\cdot\pauliVec_{v'})$.

Fix $\gamma\in(0,\pi]$.  Draw the generator axes independently and uniformly
on $\sS^2$, the gate angles from an absolutely continuous law supported on
$[-\gamma,\gamma]$ with positive density on its interior, and an independent
uniform permutation $\tau$ of the $J_n$ matching layers.  Writing
\[
    \blockV:=\prod_{v\in\gV_n}
      \exp\!\left(-\frac{\ii}{2}\theta_v\mG_v(\va_v)\right),
    \qquad
    \blockE_j:=\prod_{e\in\gM_{n,j}}
      \exp\!\left(-\frac{\ii}{2}\theta_e\mG_e(\va_e,\va'_e)\right),
\]
the mixing unitary is
\begin{equation}
    \reservoirUnitary_n(\rpar)
    :=\blockE_{\tau(J_n)}\cdots\blockE_{\tau(1)}\blockV,
    \qquad \rpar=(\skeleton,\gateangles)\sim\mathbb P_n^\omega.
    \label{eq:balanced-circuit-definition}
\end{equation}
For every fixed skeleton, the zero-angle point gives
$\reservoirUnitary_n=\identity$, and the unitary is real analytic in the gate
angles.  Its logical depth is $1+J_n\leq4$, so circuit depth is fixed by the
family rather than treated as another model hyperparameter.

\subsection{Contractive memory branches and frozen designs}
\label{sec:memory-channels}

Set $\lambda_0:=e^{-1}$.  The architecture chooses
$\lambda_+\in(\lambda_0,1)$, and each branch draws its contraction rate from
\begin{equation}
    \frac{d\nu_{\lambda_+}}{d\lambda}(\lambda)
    =\frac{\1_{(\lambda_0,\lambda_+)}(\lambda)}
      {\lambda\log(\lambda_+/\lambda_0)}.
    \label{eq:log-uniform-memory-law}
\end{equation}
Thus equal multiplicative rate intervals receive equal probability mass and
$\lambda_+$ alone sets the longest available memory scale.  The branchwise
randomization has a representational purpose beyond enforcing contraction:
a state difference propagated through branch $r$ for $\ell$ additional steps
is attenuated by $\lambda_r^\ell$, so the natural e-folding horizon is
$(-\log\lambda_r)^{-1}$.  Drawing distinct rates across the $R$ parallel
branches therefore creates a multiscale memory bank, from rapidly forgetting
to more persistent dynamics.  This heterogeneity increases the chance that a
frozen multiplex contains temporal responses compatible with the fading-memory
profile of the target functional; it is a design motivation rather than a
finite-pool performance guarantee.

Let $\rhoreset:=\lvert+\rangle\!\langle+\rvert^{\otimes n}$, with
$\lvert+\rangle=(\lvert0\rangle+\lvert1\rangle)/\sqrt2$, and define the
replacement channel $\replacementChannel(\mQ):=\Tr[\mQ]\rhoreset$.  For
$r=1,\ldots,R$, draw independently
$(\lambda_r,\rpar_r)\sim\nu_{\lambda_+}\otimes\mathbb P_n^\omega$ and set
\[
    \unitaryChannel_{r,x}(\mQ)
    :=\reservoirUnitary_n(\rpar_r)\inputUnitary_n(x)\mQ
      \inputUnitary_n(x)^\dagger\reservoirUnitary_n(\rpar_r)^\dagger.
\]
The driven reservoir channel is
\begin{equation}
    \reservoirChannel_{r,x}(\mQ)
    :=\lambda_r\unitaryChannel_{r,x}(\mQ)
      +(1-\lambda_r)\replacementChannel(\mQ),
    \label{eq:reservoir-channel}
\end{equation}
and is strictly contractive because $\lambda_r<1$.  Equivalently, this is an
input-encoding channel followed by the fixed contraction
$\mQ\mapsto\lambda_r\reservoirUnitary_n(\rpar_r)\mQ
\reservoirUnitary_n(\rpar_r)^\dagger+(1-\lambda_r)\replacementChannel(\mQ)$.
This follows the contracted-encoded QRC design principle studied by
\citet{MartinezPenaOrtega2025}, where strict contraction supplies the echo-state
and fading-memory mechanism while input dependence is carried by the encoding;
their work further analyzes when such filters distinguish input sequences.
QuaRK uses this mechanism branchwise and adds topology-aware random mixing and
the multiscale spatial multiplex described above.

For the architecture $\arch$, the frozen-design law is
\begin{equation}
    \designLaw_{\arch}
    :=\left(\nu_{\lambda_+}\otimes\mathbb P_n^\omega\right)^{\otimes R},
    \qquad \design\sim\designLaw_{\arch}.
    \label{eq:design-law}
\end{equation}
Here $\design=((\lambda_r,\rpar_r))_{r=1}^R$ records the branch-specific
parameters; the already-drawn $\projection_n$ is shared across all branches.
Once sampled, no quantum parameter is trained.

\subsection{Topology-local observables and grouped measurements}
\label{sec:topology-observables}

The complete observable bank contains all one-qubit Paulis and all two-qubit
Pauli products on graph edges:
\begin{equation}
\begin{aligned}
    \gO_V(n)&:=\{\pauliX_v,\pauliY_v,\pauliZ_v:v\in\gV_n\},\\
    \gO_E(n)&:=\{\mP_u\mP'_v:\{u,v\}\in\gE_n,
       \ \mP,\mP'\in\{\pauliX,\pauliY,\pauliZ\}\},\\
    \gO_n&:=\gO_V(n)\cup\gO_E(n),
    \qquad q_n:=|\gO_n|=12n.
\end{aligned}
    \label{eq:topology-observable-bank}
\end{equation}

Let $\gS_n=\{s_k:k=1,\ldots,K_n\}$ be a Pauli-setting cover and let
$a_n(O)$ count the settings compatible with $O$.  For the cycle, the
nine-setting construction of \cref{app:grouped-pauli-measurements} gives
$K_n=9$, with $a_n(O)=3$ for $O\in\gO_V(n)$ and $a_n(O)=1$ for
$O\in\gO_E(n)$.  We use a reusable-$n$-qubit hardware convention for the
acquisition count: repeating every setting $M$ times for each branch estimates
the entire bank and costs $RK_nM=9RM$ branch-circuit executions/state
preparations per input window.

This count measures independent circuit executions/state preparations, not
the depth or width of one execution.  Each execution propagates one $n$-qubit
branch through the $w$ inputs in the window.  Under the declared parallel gate
schedule, the per-step input injection, contraction/reset operation, and
balanced mixer have constant logical depth, so the full trajectory depth is
$O(w)$ and the quantum width is $O(n)$ under this branch-by-branch convention.
If all $R$ branches are instead physically instantiated at once, the width is
$O(Rn)$ and up to the factor $R$ can be traded from wall-clock circuit
executions into parallel hardware.  The paper reports the reusable-hardware
run count so that $R$ has one unambiguous acquisition interpretation.

\subsection{What the mixer contributes: operator spreading and direct interactions}
\label{sec:mixer-expressivity}

The universality argument later in the paper deliberately uses the identity
support point of the random circuit law.  That certificate is useful for
showing that the family has no asymptotic approximation floor, but by itself it
does not explain why the nontrivial mixer in \Eqref{eq:balanced-circuit-definition}
should be retained.  The finite-resource role of the mixer is different: it
spreads local observables across the interaction graph before measurement and
thereby exposes multiqubit functions of the encoded input directly in the
frozen feature vector.

For a Pauli string $P$, write $\operatorname{supp}(P)$ for its nonidentity
sites and $\wt(P):=|\operatorname{supp}(P)|$.  If a Hermitian operator
$A=\sum_P a_PP$ is expanded in the $n$-qubit Pauli basis, define its Pauli
interaction order by
\begin{equation}
    \pauliOrder(A):=\max\{\wt(P):a_P\neq0\}.
    \label{eq:pauli-interaction-order}
\end{equation}
For $A\subseteq\gV_n$, let
$\mathcal N_{C_n}^{\ell}(A):=\{v:\operatorname{dist}_{C_n}(v,A)\leq\ell\}$.

\begin{proposition}[Mixer-free locality]
\label{prop:identity-mixer-locality}
Set $\reservoirUnitary_n=\identity$.  If a measured Pauli observable $O$ is
supported on $A\subseteq\gV_n$, then its finite-window and causal reservoir
coordinates depend on the input history only through the projected scalar
streams
\[
    \bigl((\projection_nx_t)_j:t\leq0,\ j\in A\bigr).
\]
Consequently, vertex observables expose only one projected stream directly,
and the full bank $\gO_n$ exposes at most two projected streams directly in any
single raw feature coordinate.
\end{proposition}

\begin{theorem}[Graph-local operator spreading]
\label{thm:graph-local-operator-spreading}
Let $O$ be a Pauli string supported on $A$.  For every realization of the
balanced circuit, every Pauli component of
$\reservoirUnitary_n(\rpar)^\dagger O\reservoirUnitary_n(\rpar)$ is supported
inside $\mathcal N_{C_n}^{J_n}(A)$.  After $s$ repetitions of the same mixing
block, the support is contained in $\mathcal N_{C_n}^{sJ_n}(A)$.

Moreover, the declared continuous random-circuit law activates nonlocal
components generically.  For every oriented cycle edge $u\to v$, the
coefficient of $\pauliZ_u\pauliZ_v$ in
$\reservoirUnitary_n^\dagger\pauliX_u\reservoirUnitary_n$ is nonzero with
probability one.  For every length-two path $u-v-w$ with distinct vertices,
the coefficient of $\pauliX_u\pauliZ_v\pauliZ_w$ in
$\reservoirUnitary_n^\dagger(\pauliX_u\pauliX_v)\reservoirUnitary_n$ is
nonzero with probability one.
\end{theorem}

\begin{corollary}[Direct interaction enrichment]
\label{cor:mixer-direct-interactions}
Let
$\alpha_j(x):=\vartheta((\projection_nx)_j)$.  For a one-step branch started
from $\rhoreset$, a generic QuaRK mixer gives the $\pauliX_u$ feature a
nonzero term proportional to
$\sin\alpha_u(x)\sin\alpha_v(x)$ for every neighbor $v$ of $u$.  Likewise,
an edge feature $\pauliX_u\pauliX_v$ has, for either neighboring site $w$
outside that edge, a nonzero term proportional to
$\cos\alpha_u(x)\sin\alpha_v(x)\sin\alpha_w(x)$.
Thus the edge bank seeds genuine three-stream interactions after one mixing
block, whereas the identity mixer is restricted to direct interaction order at
most two.
\end{corollary}

\paragraph{Interpretation.}
The theorem does not claim that entanglement automatically lowers prediction
risk, nor that mixing creates frequencies absent from the encoding.  Rather,
the local encoding fixes the elementary one-coordinate trigonometric factors,
while the mixer changes which multivariate coefficient directions are exposed
to the readout.  This distinction is consistent with Fourier-based QML and QRC
expressivity analyses, where the encoding determines the accessible frequency
structure while the reservoir dynamics and readouts determine how much of that
structure is realized \citep{SchuldSwekeMeyer2021,SchuetteEtAl2026}.  It is
also aligned with recent Krylov-observability work, which relates the span of
time-evolved observables and the choice of measured operators to information
processing capacity \citep{CindrakLudgeJaurigue2025,CindrakJaurigueLudge2025}.

The cycle now has a precise resource role.  It is connected, so it imposes no
permanent decomposition into isolated components, while its bounded degree
permits every entangling block to be scheduled in $J_n\leq3$ matching layers.
The graph controls the light cone in which interactions can be generated; the
edge observables provide two-site seeds and therefore expose one higher direct
interaction order than vertex observables after a generic block.  These are
finite-resource representation statements.  They do not contradict the later
identity-support universality proof, whose universal nonlinear readout can
compose simpler coordinates given sufficiently many resources.

\subsection{Exact and measured feature maps}

For a finite window $\vx=x_{-w+1:0}$, the state of branch $r$ is
\[
    \rho_r^{\design}(\vx)
    :=\reservoirChannel_{r,x_0}\circ\cdots\circ
      \reservoirChannel_{r,x_{-w+1}}(\rhoreset).
\]
Strict contraction gives the corresponding causal state
$\rho_{r,\infty}^{\design}(\vx^-)$ by moving the initialization to the remote
past.  The branch feature block is
$\phi_r^{\design}(\vx):=(\Tr[O\rho_r^{\design}(\vx)])_{O\in\gO_n}$.
With $m_{\arch}:=Rq_n$, the exact finite-window feature map is
\begin{equation}
    \featureMap_{\design}(\vx)
    :=\bigl(\phi_1^{\design}(\vx),\ldots,\phi_R^{\design}(\vx)\bigr)
    \in[-1,1]^{m_{\arch}},
    \label{eq:exact-window-features}
\end{equation}
and $\featureMap_{\design,\infty}$ is defined analogously from the causal
states.  We write $\measuredFeatureMap_{\design}(\vx;\chi)$ for the estimator
obtained from the grouped measurement record $\chi$.  Exact and measured
features remain distinct throughout the analysis.

\subsection{Mat\'ern readout and risks}
\label{sec:readout-risks}

Feature dimension changes with the architecture, so an unnormalized Euclidean
kernel distance would mechanically grow when more bounded coordinates are
concatenated.  QuaRK removes that nuisance at the model-definition level.  For
$z,z'\in\sR^p$, define the root-mean-square feature distance
\begin{equation}
    \rmsMetric_p(z,z')
    :=\frac{\norm{z-z'}_2}{\sqrt p}.
    \label{eq:rms-feature-metric}
\end{equation}
Let $\kernel_p$ be the normalized Mat\'ern kernel with smoothness $\nu>1$ and
length scale $\xi>0$ evaluated at $\rmsMetric_p(z,z')$.  To remove any
implementation ambiguity, throughout the paper ``Mat\'ern$(\nu,\xi)$'' means
\begin{equation}
    \kernel_p(z,z')
    :=\frac{2^{1-\nu}}{\Gamma(\nu)}
      \left(\frac{\sqrt{2\nu}\,\rmsMetric_p(z,z')}{\xi}\right)^{\nu}
      K_{\nu}\!\left(\frac{\sqrt{2\nu}\,\rmsMetric_p(z,z')}{\xi}\right),
    \qquad \kernel_p(z,z):=1,
    \label{eq:matern-kernel-definition}
\end{equation}
where $K_{\nu}$ is the modified Bessel function of the second kind, and let
$\RKHS_p$ be its RKHS.  Equivalently, this is the ordinary fixed-$(\nu,\xi)$
Mat\'ern kernel applied to $z/\sqrt p$.  The $1/\sqrt p$ factor is fixed once
by the method, is label independent, and is never retuned across
architectures.  For every finite $p$ the rescaling is invertible, so it does
not remove information or alter Mat\'ern universality on compact feature
images.

The Ivanov ball is
$\RKHS_p^\Lambda:=\{h\in\RKHS_p:\norm{h}_{\RKHS_p}\leq\Lambda\}$.  Kernel
normalization gives $|h(z)|\leq\Lambda$, while the section map is Lipschitz
in the RMS metric with the dimension-free constant
\[
    \kLip:=\frac1\xi\sqrt{\frac\nu{\nu-1}}.
\]
Equivalently, its Lipschitz constant with respect to the raw Euclidean feature
vector is $\kLip/\sqrt p$.  This normalization prevents feature count alone
from being interpreted as either architectural improvement or degradation;
the empirical protocol keeps one fixed Mat\'ern specification and checks the
resulting similarity geometry directly.
The radius $\Lambda$ is fixed before fitting; \Secref{sec:dependent-generalization}
provides a direct calibration from a requested generalization tolerance.

For a frozen design $\design$, define the causal, finite-window, and
operational risks together by
\begin{align*}
    \risk_{\design,\infty}(h)
    &:=\E\!\left[(h(\featureMap_{\design,\infty}(X_{-\infty:0}))-Y_0)^2\right],\\
    \risk_{\design,w}(h)
    &:=\E\!\left[(h(\featureMap_{\design}(X_{-w+1:0}))-Y_0)^2\right],\\
    \oprisk_{\design,w}(h)
    &:=\E\!\left[(h(\measuredFeatureMap_{\design}(X_{-w+1:0}))-Y_0)^2\right],
\end{align*}
where the last expectation also averages over fresh measurement outcomes.
For nonempty $\mathcal I\subseteq[N]$, the measured empirical risk is
\[
    \erisk_{\design,w}^{\mathcal I}(h)
    :=\frac1{|\mathcal I|}\sum_{j\in\mathcal I}
      \bigl(h(\measuredFeatureMap_{\design}(X_{j\Delta-w+1:j\Delta}))
      -Y_{j\Delta}\bigr)^2,
    \qquad
    \erisk_{\design,w}:=\erisk_{\design,w}^{[N]}.
\]
Finally, make the readout-radius dependence of the causal approximation error
explicit:
\begin{equation}
    \epsApp(\design;\Lambda)
    :=\inf_{h\in\RKHS_{m_{\arch}}^\Lambda}
      \bigl\{\risk_{\design,\infty}(h)-\risk^\star\bigr\}.
    \label{eq:design-approximation-error}
\end{equation}
This quantity links the width-resolved representation analysis in
\Secref{sec:width-resolved-approximation} to the finite-pool performance theorem.
